% Draft scalability section. NOT auto-inserted into the paper.
% Generated from results/scalability/optimized_segmental_v1/ at commit 564995efd056d7d33984f0ca1532386e6140ea0c.
% Every number below is read from timing_summary.csv / memory_summary.csv /
% complexity_fits.json. Regenerate with scripts/scalability/report.py.

\section{Computational scalability}
\label{sec:scalability}

Exact inference in the segmental partial-order skill model is a semi-Markov
forward--backward recursion over candidate blocks. With the maximum segment width
$D$ bounded, the factorised forward pass costs $O(N[JK^2 + JDK])$ per all-trace
sweep, against $O(NJDK^2)$ for the direct recursion. We measured the optimized
backend against the frozen reference engine on a grid of small problems before any
scaling measurement: the worst absolute discrepancy in $\alpha$ and in $\log Z$ was
1.1e-13, against a
pre-registered tolerance of $10^{-10}$, with identical $-\infty$ support patterns
and bit-identical candidate score tables. The two implementations compute the same
numbers.

\paragraph{Protocol.}
Every configuration ran in its own process on a single pinned thread, with all BLAS
and OpenMP thread counts set to one and configurations executed strictly
sequentially. Each point reports the median of at least fifteen timed repetitions
after three untimed warm-ups, with repetitions continued until the bootstrap
$95\%$ interval for the median had relative half-width at most $5\%$. Expensive
points are allowed five repetitions and are marked as such. Wall-clock and process
CPU time were both recorded and agree throughout.

\paragraph{Trace length.}
With $D$ bounded, runtime grows approximately linearly in trace length. Over
$J \in [24$-$1024]$ the optimized forward pass scales as $J^{1.38}$ and a
complete plain sweep as $J^{0.99}$ (Table~\ref{tab:scal-J}).

\begin{table}[t]
  \centering
  \small
  \begin{tabular}{rrrrrr}
    \toprule
    $J$ & legal blocks & forward & FFBS & plain sweep & peak RSS \\
    \midrule
    1024 & 162,800 & 319.8\,ms & 478.9\,ms & 528.8\,ms & 4.78\,GiB \\
    192 & 29,680 & 13.7\,ms & 42.7\,ms & 61.8\,ms & 559\,MiB \\
    24 & 2,800 & 1.5\,ms & 4.6\,ms & 11.9\,ms & 126\,MiB \\
    384 & 60,400 & 63.6\,ms & 122.5\,ms & 147.1\,ms & 1.42\,GiB \\
    48 & 6,640 & 3.1\,ms & 9.7\,ms & 20.4\,ms & 164\,MiB \\
    768 & 121,840 & 192.7\,ms & 309.6\,ms & 350.3\,ms & 4.12\,GiB \\
    96 & 14,320 & 6.4\,ms & 20.4\,ms & 36.0\,ms & 263\,MiB \\
    \bottomrule
  \end{tabular}
  \caption{Trace-length scaling at $N=16$, $K=10$, $A=20$, $D\in[3,12]$. Medians
  over at least fifteen timed repetitions in a dedicated process.}
  \label{tab:scal-J}
\end{table}

\paragraph{Skills, corpus, width and roles.}
A complete plain sweep scales as $K^{0.79}$ in the size of the skill
library and $N^{0.72}$ in corpus size. Widening the maximum segment width
increases cost through the number of legal candidate blocks rather than through $D$
directly, and we report it against that count. The size of the canonical-action
inventory $A$ acts through the density of the induced role graph, not through $A$
alone: with every skill supported on all $A$ roles the candidate-table construction
grows sharply, while with a fixed ten-role support drawn from a size-$A$ vocabulary
it is close to flat. We report the two regimes separately and never average them.

\paragraph{Marginalisation.}
A sweep that performs no structural update executes identical code in both arms, so
marginalising the partial order is free on that path; the entire overhead of
FULL-MARG falls on the scheduled structural update, and at the registered cadence of
one structural sweep in ten it amortises to the ratio reported in the supplement.

\paragraph{Operating point.}
At the scale we anticipate for real data --- $N=100$ traces of length $J=200$ over
$K=20$ skills and $A=50$ canonical actions, with a ten-role support per skill and
$D\in[3,12]$ --- a plain sweep takes 648.3\,ms and the process peaks at
3.66\,GiB resident. These are throughput measurements; they bear on cost per
sweep and not on the number of sweeps required.

\paragraph{Memory.}
The binding constraint on trace length is memory rather than time. The current
implementation stores a dense $(J, J{+}1, K)$ candidate score table per trace, so
score-table memory is $O(NJ^2K)$. Storing only the $D_{\max} - D_{\min} + 1$
legal durations per start would reduce this to $O(NJDK)$, but
\emph{this layout is not implemented here}: every banded figure we quote is
arithmetic on array shapes rather than a measurement.
